\section*{M3: $k=3$ branching topology}

\textbf{Construction:} root joins two independent 2-leaf binary nodes,
$T = \mu(\mu(L_0,L_1), \mu(L_2,L_3))$ ($u_1,u_2 \to \text{root}$, mission's
named branching shape), same homogeneous gated-planar-rotation law and
conventions as the chain case.

\textbf{Theorem (new this pass).}
\[
  E_T^{\mathrm{proj}}(\eta) \;=\; \eta^2\sqrt{1-\eta^2}, \qquad \eta\in[0,1].
\]
Proved the same way as the chain case (exact symbolic substitution
through the repository's own evaluator; $\|\text{diff}\|^2 =
\eta^4(1-\eta^2)$ identically).

\textbf{Ratio to the universal bound} ($2\eta$, same as chain, since both
have $k=3$):
\[
  \frac{E_T^{\mathrm{proj}}(\eta)}{2\eta} = \frac12\,\eta\sqrt{1-\eta^2},
\]
maximized at the same $\eta^\star=1/\sqrt2$, with maximum value
$\boxed{1/4}$ (relative gap $3/4$) — branching is a substantially
\emph{less} efficient construction than chain for this law family:
$1/4$ vs.\ chain's $3/4$ best-ratio, a factor of 3 apart, for the exact
same universal bound. This shows the $k=3$ certified-gap floor is
\textbf{topology-dependent} even within one fixed law family and fixed
$(n,r)$ — a genuinely new, exact, closed-form confirmation of what the
prior session's terminal-status document could only describe
qualitatively ("consistent with the single-topology gated-rotation
construction not generalizing cleanly across topologies").

\textbf{Status:} neither topology saturates the $k=3$ universal bound at
any $\eta$; chain is the better of the two named topologies for this law
family, by an exact factor of 3 at the shared optimal $\eta^\star=1/\sqrt2$.
